\chapter{Feline fibrations and derived constructions}\label{ch:fibrations}

\begin{assumption}\label{ass:ch10}
Throughout this chapter $\cC$ has an initial object $0$ and a terminal object
$1$. Additional hypotheses are stated in each section.
\end{assumption}

\section{The leg fibration}\label{sec:leg-fibration}

\begin{definition}\label{def:reindexing-functor}
For a function $\varphi\colon A\to B$ and a cat $K'$ with leg set $B$, the
\emph{reindexing} of $K'$ along $\varphi$ is the cat
$\varphi^*K'\defeq\langle L'_{\varphi(\alpha)}\rangle_{\alpha\in A}$, and the
\emph{canonical lift} of $\varphi$ is the morphism
$\bar\varphi=(\varphi,(\id_{L'_{\varphi(\alpha)}})_\alpha)\colon\varphi^*K'\to K'$.
\end{definition}

\begin{theorem}\label{thm:fibration}
The leg functor $\legs\colon\Fel(\cC)\to\Set$ is a Grothendieck fibration. A
feline morphism is cartesian if and only if it is reindexing in the sense of
\cref{def:reindexing}, that is, all of its leg components are isomorphisms.
The canonical lifts form a cleavage which, on normalized cats, is split:
$(\psi\varphi)^*=\varphi^*\psi^*$ and $\id^*=\id$. The fibre over a set $A$ is
equivalent to $\cC^A$.
\end{theorem}

\begin{proof}
Let $(\varphi,F)\colon K\to K'$ be reindexing. Given $(\psi,G)\colon M\to K'$
and $\chi\colon\legs M\to A$ with $\varphi\chi=\psi$, a morphism
$(\chi,H)\colon M\to K$ satisfies $(\varphi,F)(\chi,H)=(\psi,G)$ if and only if
$F_{\chi(c)}H_c=G_c$ for all $c$, that is, $H_c=F_{\chi(c)}^{-1}G_c$. So
$(\varphi,F)$ is cartesian. In particular canonical lifts are cartesian, and
$\legs$ is a fibration. Conversely, suppose $(\varphi,F)\colon K\to K'$ is
cartesian. Since $\bar\varphi\colon\varphi^*K'\to K'$ is also cartesian over
$\varphi$, the universal properties provide vertical morphisms
$(\id,V)\colon K\to\varphi^*K'$ and $(\id,W)\colon\varphi^*K'\to K$ with
$\bar\varphi(\id,V)=(\varphi,F)$ and $(\varphi,F)(\id,W)=\bar\varphi$; uniqueness
in the universal properties shows that they are mutually inverse. Comparing
components, $V_\alpha=F_\alpha$, and $F_\alpha$ is an isomorphism. The
splitting identities hold on normalized cats by inspection, and the statement
about fibres is \cref{prop:fibre-legs}.
\end{proof}

\begin{remark}\label{rem:chloe-fibration}
By \cref{prop:chloe-basic}(iii), cats satisfying the Chloe Condition are closed
under reindexing, so the canonical lifts of \cref{thm:fibration} remain in the
full subcategory $\Fel^\Chl(\cC)$, and the restriction of $\legs$ to
$\Fel^\Chl(\cC)$ is again a fibration. The opcartesian lifts of
\cref{thm:opfibration}, when they exist, need not remain in $\Fel^\Chl(\cC)$: in
$\Set$, the opcartesian lift of $2\to1$ at the Chloe cat $\langle1,1\rangle$ has
the two-element set as its leg, which is not connected.
\end{remark}


\begin{proposition}\label{prop:groth}
Let $\cC^{(-)}\colon\Set^{\op}\to\CAT$ be the functor $A\mapsto\cC^A$,
$\varphi\mapsto\varphi^*$ (precomposition). Then $\Fel^{\nrm}(\cC)$ is isomorphic,
over $\Set$, to the Grothendieck construction $\El\cC^{(-)}$ of
\cref{def:groth-construction}. Hence $\Fel(\cC)$ is equivalent over $\Set$ to
$\El\cC^{(-)}$.
\end{proposition}

\begin{proof}
An object of $\El\cC^{(-)}$ is a pair $(A,L)$ with $L\in\cC^A$, and a morphism
$(A,L)\to(B,L')$ is a pair $(\varphi,F)$ with $F\colon L\to\varphi^*L'$ in $\cC^A$,
that is, a family $F_\alpha\colon L_\alpha\to L'_{\varphi(\alpha)}$. The composite
$(\psi,G)\circ(\varphi,F)=(\psi\varphi,\varphi^*(G)\circ F)$ has components
$G_{\varphi(\alpha)}F_\alpha$. This is exactly $\Fel^{\nrm}(\cC)$ under
\cref{prop:automatic}; the second assertion follows from
\cref{prop:normalization}.
\end{proof}

\begin{theorem}\label{thm:opfibration}
The leg functor is an opfibration, and hence a bifibration, if and only if
$\cC$ has small coproducts. In that case a morphism $(\varphi,F)\colon K\to K'$
is opcartesian if and only if, for each $\beta\in\legs K'$, the family
$(F_\alpha\colon L_\alpha\to L'_\beta)_{\alpha\in\varphi^{-1}(\beta)}$ is a
coproduct cocone. In particular the opcartesian lift of $\varphi\colon A\to B$
at $K$ is $(\varphi,\inj)\colon K\to\varphi_!K\defeq\langle\coprod_{\varphi(\alpha)=\beta}L_\alpha\rangle_{\beta\in B}$.
\end{theorem}

\begin{proof}
Suppose $(\varphi,F)$ has the stated property. Given $(\psi,G)\colon K\to M$
and $\chi\colon B\to\legs M$ with $\chi\varphi=\psi$, a morphism
$(\chi,H)\colon K'\to M$ with $(\chi,H)(\varphi,F)=(\psi,G)$ must satisfy
$H_\beta F_\alpha=G_\alpha$ for all $\alpha\in\varphi^{-1}(\beta)$; since
$(F_\alpha)_{\alpha\in\varphi^{-1}(\beta)}$ is a coproduct cocone and
$G_\alpha\colon L_\alpha\to M_{\psi(\alpha)}=M_{\chi(\beta)}$, there is exactly
one such $H_\beta$. So $(\varphi,F)$ is opcartesian, and if $\cC$ has small
coproducts the morphisms $(\varphi,\inj)$ are opcartesian lifts.

Conversely, let $(\varphi,F)\colon K\to K'$ be opcartesian and $\beta\in B$.
Given a cocone $(g_\alpha\colon L_\alpha\to Z)_{\alpha\in\varphi^{-1}(\beta)}$,
let $M$ be obtained from $K'$ by replacing the leg $\beta$ by $Z$, and let
$(\varphi,G)\colon K\to M$ have components $g_\alpha$ for
$\alpha\in\varphi^{-1}(\beta)$ and $F_\alpha$ otherwise. Taking $\chi=\id_B$,
opcartesianness yields a unique vertical morphism $(\id,H)\colon K'\to M$ with
$(\id,H)(\varphi,F)=(\varphi,G)$; its component $h_0=H_\beta$ satisfies
$h_0F_\alpha=g_\alpha$ for all $\alpha\in\varphi^{-1}(\beta)$. If
$h\colon L'_\beta\to Z$ is any morphism with the same property, let $(\id,H^h)$
be the vertical morphism with component $h$ at $\beta$ and component $H_{\beta'}$
at every $\beta'\neq\beta$. Then $(\id,H^h)(\varphi,F)=(\varphi,G)$, so
$(\id,H^h)=(\id,H)$ and $h=h_0$. Hence $(F_\alpha)_{\alpha\in\varphi^{-1}(\beta)}$
is a coproduct cocone. Finally, if $\legs$ is an opfibration, let $(X_s)_{s\in S}$
be a small family in $\cC$ and take an opcartesian lift of $S\to1$ at
$\langle X_s\rangle_S$; by what has been shown, its unique leg is a coproduct of
the $X_s$.
\end{proof}

\begin{proposition}[Factorizations]\label{prop:factorization}
Every feline morphism $(\varphi,F)\colon K\to K'$ factors as a vertical morphism
followed by a cartesian one,
\[
(\varphi,F)=\bar\varphi\circ(\id_A,F)\colon K\too\varphi^*K'\too K',
\]
and this factorization is unique up to a unique vertical isomorphism. If $\cC$
has small coproducts it also factors as an opcartesian morphism followed by a
vertical one, $(\varphi,F)=(\id_B,\tilde F)\circ(\varphi,\inj)$, where
$\tilde F_\beta\colon\coprod_{\varphi(\alpha)=\beta}L_\alpha\to L'_\beta$ is induced by
the $F_\alpha$.
\end{proposition}

\begin{proof}
Both identities are verified on gaits and components. Uniqueness up to
unique vertical isomorphism is the general uniqueness of cartesian (resp.\
opcartesian) lifts.
\end{proof}

\section{Tail-absorbing morphisms}\label{sec:tail-absorbing}

\Cref{rem:morphism-functorial} described feline morphisms as pairs $(u,\theta)$
with $u\colon\Fsh_A\to\Fsh_B$ leg-preserving. If $u$ is allowed to send legs to
the tail, we obtain a larger class of morphisms.

\begin{definition}\label{def:tail-absorbing}
Let $K$ and $K'$ be cats with leg sets $A$ and $B$, and write
$B_{\mathsf t}=B\sqcup\{\mathsf t\}$. A \emph{tail-absorbing morphism}
$K\rightsquigarrow K'$ is a pair $(\varphi,F)$ where $\varphi\colon A\to B_{\mathsf t}$ and
$F=(F_\alpha)_{\alpha\in\varphi^{-1}(B)}$ with $F_\alpha\colon L_\alpha\to L'_{\varphi(\alpha)}$.
The composite of $(\varphi,F)\colon K\rightsquigarrow K'$ and
$(\psi,G)\colon K'\rightsquigarrow K''$ has gait
$\alpha\mapsto\mathsf t$ if $\varphi(\alpha)=\mathsf t$ and $\alpha\mapsto\psi(\varphi(\alpha))$
otherwise, and components $G_{\varphi(\alpha)}F_\alpha$ for those $\alpha$ with
$\psi(\varphi(\alpha))\neq\mathsf t$. The resulting category is
denoted $\Fel^{\ta}(\cC)$.
\end{definition}

A leg sent to $\mathsf t$ is \emph{absorbed into the tail}; its component would
be the unique morphism $L_\alpha\to T'$, and is omitted.

\begin{proposition}\label{prop:tail-absorbing-functorial}
Tail-absorbing morphisms $K\rightsquigarrow K'$ correspond bijectively to pairs
$(u,\theta)$ where $u\colon\Fsh_A\to\Fsh_B$ is a functor fixing $\mathsf h$ and
$\mathsf t$ and sending each $\ell_\alpha$ to some $\ell_\beta$ or to $\mathsf t$,
and $\theta\colon K\Rightarrow K'\circ u$ is a natural transformation; composition
corresponds to $(v,\zeta)\circ(u,\theta)=(vu,\zeta u\circ\theta)$.
\end{proposition}

\begin{proof}
As in \cref{rem:morphism-functorial}: such $u$ correspond to functions
$\varphi\colon A\to B_{\mathsf t}$, and the components of $\theta$ at $\mathsf h$,
at $\mathsf t$ and at legs sent to $\mathsf t$ are forced, all naturality
squares involving them commuting by \cref{lem:parallel}.
\end{proof}

\begin{proposition}\label{prop:kleisli}
Let $\mathbb T$ be the monad $K\mapsto K\sqcup\Caudal$ on $\Fel(\cC)$, with unit
the coprojection $K\to K\sqcup\Caudal$ and multiplication
$K\sqcup\Caudal\sqcup\Caudal\to K\sqcup\Caudal$ the codiagonal on the two copies
of $\Caudal$. Then $\Fel^{\ta}(\cC)$ is isomorphic to the Kleisli category
$\Kl(\mathbb T)$, by an isomorphism that is the identity on objects.
\end{proposition}

\begin{proof}
That $\mathbb T$ is a monad follows from the universal property of
coproducts: both sides of each monad law are morphisms out of a coproduct, and
they agree on each summand. Using the coproduct of \cref{prop:coproducts}, the leg set
of $K'\sqcup\Caudal$ is $B\sqcup\{*\}$, which we identify with $B_{\mathsf t}$,
and the leg at $\mathsf t$ is $1$. A Kleisli morphism $K\to K'\sqcup\Caudal$
therefore consists of a function $\varphi\colon A\to B_{\mathsf t}$, components
$L_\alpha\to L'_{\varphi(\alpha)}$ for $\varphi(\alpha)\in B$, and the unique
components $L_\alpha\to1$ otherwise; this is a tail-absorbing morphism. Kleisli
composition $\mu\circ\mathbb T(g)\circ f$ sends a leg absorbed by $f$ to the
tail (via the codiagonal), and otherwise applies $g$; this is the composition
of \cref{def:tail-absorbing}.
\end{proof}

\begin{example}\label{ex:partial}
For $\cC=\one$, $\Fel(\one)\iso\Set$, and $\Fel^{\ta}(\one)$ is the category of
small sets and partial functions, a leg absorbed into the tail being an
element at which the function is undefined. This category is well known to be
equivalent to the category $\Setp$ of pointed sets.
\end{example}

\begin{remark}\label{rem:head-absorbing}
One may also let $u$ send legs to the head. A component $L_\alpha\to H'$ is then
required, and such a component exists only if $L_\alpha$ admits a morphism to
an initial object; if the initial object of $\cC$ is strict, this forces
$L_\alpha$ to be initial. In $\Set$, for example, only empty legs may be absorbed
into the head. We do not pursue \emph{head-absorbing} morphisms.
\end{remark}

\section{Kan extensions along the one-legged inclusion}\label{sec:kan}

\begin{proposition}\label{prop:kan}
Let $G\colon\cC\to\cE$ be a functor.
\begin{enumerate}
\item If $\cE$ has small coproducts, the functor $\widehat G\colon\Fel(\cC)\to\cE$,
  $\widehat G(K)=\coprod_{\alpha}G(L_\alpha)$, with $\widehat G(\psi,F)$ restricting on
  the $\alpha$-th summand to $\inj_{\psi(\alpha)}\circ G(F_\alpha)$, is a left Kan
  extension of $G$ along $\iota$, and $\widehat G\circ\iota\iso G$.
\item If $\cC$ has small coproducts, $G\circ\body$ is a right Kan extension of
  $G$ along $\iota$.
\end{enumerate}
\end{proposition}

\begin{proof}
(i) We identify $\widehat G\iota X=\coprod_1G(X)$ with $G(X)$. Let
$H\colon\Fel(\cC)\to\cE$ be a functor. Given $\sigma\colon G\Rightarrow H\iota$,
let $\widehat\sigma_K\colon\widehat G(K)\to H(K)$ be the morphism whose restriction to
the $\alpha$-th summand is $H(j_\alpha)\circ\sigma_{L_\alpha}$. For
$(\psi,F)\colon K\to K'$, \eqref{eq:leg-inclusion-naturality} and naturality of
$\sigma$ give
\[
H(\psi,F)H(j_\alpha)\sigma_{L_\alpha}=H(j_{\psi(\alpha)})H(\iota F_\alpha)\sigma_{L_\alpha}
=H(j_{\psi(\alpha)})\sigma_{L'_{\psi(\alpha)}}G(F_\alpha),
\]
which is the restriction of $\widehat\sigma_{K'}\circ\widehat G(\psi,F)$ to the
$\alpha$-th summand; so $\widehat\sigma$ is natural. Conversely, a natural
transformation $\tau\colon\widehat G\Rightarrow H$ restricts to
$\tau\iota\colon G\Rightarrow H\iota$. These constructions are mutually inverse:
$\widehat\sigma\iota=\sigma$ because $j_*=\id$ for a unipedal cat; and the
restriction of $\widehat{\tau\iota}_K$ to the $\alpha$-th summand is
$H(j_\alpha)\tau_{\iota L_\alpha}=\tau_K\circ\widehat G(j_\alpha)$, by naturality of
$\tau$, which is the restriction of $\tau_K$, since $\widehat G(j_\alpha)$ is the
$\alpha$-th coprojection. The resulting bijection
$\mathrm{Nat}(\widehat G,H)\iso\mathrm{Nat}(G,H\iota)$ is natural in $H$, which is
the universal property of a left Kan extension.

(ii) Let $\eta\colon\id\Rightarrow\iota\body$ and $\varepsilon\colon\body\iota\Rightarrow\id$
be the unit and counit of $\body\dashv\iota$. For $H\colon\Fel(\cC)\to\cE$, the maps
\begin{gather*}
\mathrm{Nat}(H,G\body)\too\mathrm{Nat}(H\iota,G),\qquad \tau\mapsto G\varepsilon\circ\tau\iota,\\
\mathrm{Nat}(H\iota,G)\too\mathrm{Nat}(H,G\body),\qquad \rho\mapsto\rho\body\circ H\eta,
\end{gather*}
are mutually inverse. Indeed, by naturality of $\rho$ and a triangle identity,
\[
G\varepsilon\circ\rho\body\iota\circ H\eta\iota=\rho\circ H(\iota\varepsilon\circ\eta\iota)=\rho,
\]
and by naturality of $\tau$ and the other triangle identity,
\[
G\varepsilon\body\circ\tau\iota\body\circ H\eta=G(\varepsilon\body\circ\body\eta)\circ\tau=\tau.
\]
\end{proof}

\begin{remark}\label{rem:kan-pointwise}
The comma category $\iota\downarrow K$ is isomorphic to $\coprod_\alpha\cC/L_\alpha$,
each $\cC/L_\alpha$ having the terminal object $(L_\alpha,\id)$; by
\cref{lem:colim-terminal}, $\widehat G(K)$ is therefore also the colimit of $G$
over $\iota\downarrow K$, so the left Kan extension of (i) is pointwise. Dually,
$K\downarrow\iota$ is the category of cocones on the legs of $K$, whose initial
object is the coproduct cocone, so the right Kan extension of (ii) is pointwise
as well.
\end{remark}

\begin{remark}\label{rem:kan-honest}
Part (i) is the coproduct-preserving extension of \cref{thm:universal-property};
part (ii) is an instance of the general fact that right Kan extension along a
right adjoint is precomposition with its left adjoint.
\end{remark}

\section{The classifying space}\label{sec:nerve}

Since $\Fel(\cC)$ is not small, we fix a second universe $\mathbb U'$ with
$\mathbb U\in\mathbb U'$ and regard $\Fel(\cC)$ as a $\mathbb U'$-small category
when $\cC$ is $\mathbb U$-small. The nerve and classifying space are then
defined as usual, relative to $\mathbb U'$.

\begin{theorem}\label{thm:contractible}
Let $\cC$ be small. The classifying space of $\Fel(\cC)$ is contractible.
More generally, every full subcategory of $\Fel(\cC)$ that contains a legless cat
or a terminal cat has a contractible classifying space; this applies to
$\Kit(\cC)$, to $\Fel_1(\cC)$, and to the full subcategory of cats with at least
one leg.
\end{theorem}

\begin{proof}
A category $\cD$ with an initial object $i$ has a natural transformation from
the constant functor at $i$ to the identity functor; the induced homotopy
between the identity of $B\cD$ and a constant map shows that $B\cD$ is
contractible \cite{Segal1968,Quillen1973}. The dual argument applies to a
terminal object. A legless cat is initial, and the caudal cat is terminal, in
any full subcategory containing them.
\end{proof}

\begin{remark}\label{rem:nerve-honest}
\Cref{thm:contractible} says that the homotopy type of the category of cats
carries no information. This is a general property of categories with an
initial or terminal object and has nothing to do with cats. The groupoid of
cats and isomorphisms, considered next, is more informative.
\end{remark}

\section{The census of cats}\label{sec:census}

\begin{definition}\label{def:multiplicity}
Let $\Iso(\cC)$ be the class of isomorphism classes of objects of $\cC$. The
\emph{multiplicity function} of a cat $K$ is the function
$m_K\colon\Iso(\cC)\to\mathrm{Card}$,
$m_K([X])=\bigl|\{\alpha\in\legs K:L_\alpha\iso X\}\bigr|$.
\end{definition}

The function $m_K$ has small support, and its values are small cardinals.

\begin{theorem}[Census]\label{thm:census}
Two cats $K$ and $K'$ are isomorphic if and only if $m_K=m_{K'}$. Every function
$m\colon\Iso(\cC)\to\mathrm{Card}$ with small support and small values is the
multiplicity function of some cat. Hence the isomorphism classes of
$\cC$-cats are in bijection with such functions.
\end{theorem}

\begin{proof}
If $(\varphi,F)\colon K\to K'$ is an isomorphism then $\varphi$ is a bijection
with $L_\alpha\iso L'_{\varphi(\alpha)}$ (\cref{prop:iso}), so $\varphi$ restricts
to bijections between the sets of legs in each isomorphism class, and
$m_K=m_{K'}$. Conversely, if $m_K=m_{K'}$, choose for each class $[X]$ a
bijection between the legs of $K$ in the class $[X]$ and those of $K'$, and for
each leg an isomorphism with the corresponding leg (axiom of choice); this
defines an isomorphism $K\to K'$ by \cref{prop:iso}. Given $m$, choose a
representative $X$ of each class in the support and take the cat with $m([X])$
legs equal to $X$ for each such class.
\end{proof}

\begin{theorem}\label{thm:automorphism}
Let $K$ be a cat, and for each class $[X]$ in the support of $m_K$ fix a
representative $X$ and let $S_{[X]}$ be the set of legs of $K$ in the class
$[X]$. Then
\[
\Aut(K)\iso\prod_{[X]}\Aut(X)\wr\Sym(S_{[X]}),
\]
where for a group $G$ and a set $S$, $G\wr\Sym(S)$ denotes the set
$G^S\times\Sym(S)$ with multiplication $(f,\sigma)(g,\tau)=((f\circ\tau)\cdot g,\sigma\tau)$,
the product $\cdot$ being pointwise.
\end{theorem}

\begin{proof}
Choosing isomorphisms $L_\alpha\iso X$ for each leg in the class $[X]$, we may
assume, by transport of structure, that every leg of $K$ in the class $[X]$
equals $X$. An automorphism $(\varphi,F)$ of $K$ has bijective gait preserving
isomorphism classes of legs, so $\varphi$ is a family
$(\varphi_{[X]})\in\prod_{[X]}\Sym(S_{[X]})$, and its components form a family
$(F_\alpha)\in\prod_{[X]}\Aut(X)^{S_{[X]}}$. Composition is
$(\varphi,F)(\psi,G)=(\varphi\psi,(F_{\psi(\alpha)}G_\alpha)_\alpha)$, which, on each
factor, is the multiplication of the wreath product with $f=F$, $g=G$,
$\sigma=\varphi$ and $\tau=\psi$.
\end{proof}

\begin{example}\label{ex:automorphisms}
\begin{enumerate}
\item In $\Fel(\one)\iso\Set$, $\Aut(A)=\Sym(A)$.
\item In $\Fel(\Set)$, a cat with $n$ legs each of which is a singleton has
  automorphism group $\Sym(n)$, and a cat with $n$ legs each of which has two
  elements has automorphism group $(\mathbb Z/2)\wr\Sym(n)$, the hyperoctahedral
  group of order $2^nn!$.
\item A quadrupedal $\Set$-cat whose legs have $1$, $1$, $2$ and $3$ elements has
  automorphism group $\Sym(2)\times(\mathbb Z/2)\times\Sym(3)$, of order $24$.
\end{enumerate}
\end{example}

\section{The kitten rig}\label{sec:kitten-rig}

In this section $\cC$ is essentially small and symmetric monoidal. Let
$\Iso(\cC)$ be the set of isomorphism classes of objects, which is a commutative
monoid under $[X][Y]=[X\otimes Y]$ with unit $[I]$. For a commutative monoid
$M$ we write $\mathbb N[M]$ for the monoid semiring: finitely supported functions
$M\to\mathbb N$ with pointwise addition and convolution product.

\begin{theorem}\label{thm:kitten-rig}
The set $\mathcal R(\cC)$ of isomorphism classes of $\cC$-kittens is a
commutative semiring under $[K]+[K']=[K\sqcup K']$ and $[K][K']=[K\otimes K']$,
with zero $[\Legless]$ and unit $[\mathbf I]$. The map $[K]\mapsto m_K$ is an
isomorphism of semirings $\mathcal R(\cC)\iso\mathbb N[\Iso(\cC)]$, and the leg
count $[K]\mapsto|\legs K|$ corresponds to the augmentation
$\mathbb N[\Iso(\cC)]\to\mathbb N$, $m\mapsto\sum_xm(x)$.
\end{theorem}

\begin{proof}
The operations are well defined on isomorphism classes since $\sqcup$ and
$\otimes$ are functors. By \cref{thm:census}, $[K]\mapsto m_K$ is a bijection
onto finitely supported functions $\Iso(\cC)\to\mathbb N$. It is additive since
$m_{K\sqcup K'}=m_K+m_{K'}$ (\cref{prop:coproducts}), and multiplicative since
the legs of $K\otimes K'$ are the $L_\alpha\otimes L'_\beta$, so that
$m_{K\otimes K'}(z)=\sum_{xy=z}m_K(x)m_{K'}(y)$. The semiring axioms for
$\mathcal R(\cC)$ follow, or may be read off from \cref{thm:monoidal} and
\cref{prop:tensor-distributes}. The leg count is $\sum_xm_K(x)$.
\end{proof}

\begin{example}\label{ex:kitten-rig}
\begin{enumerate}
\item $\mathcal R(\one)\iso\mathbb N$.
\item Let $\cC$ be the category of finite sets with the cocartesian monoidal
  structure $\sqcup$. Then $\Iso(\cC)\iso(\mathbb N,+)$ and
  $\mathcal R(\cC)\iso\mathbb N[x]$, the semiring of polynomials with
  natural-number coefficients: a kitten whose legs have $n_1,\dots,n_k$
  elements corresponds to $x^{n_1}+\dots+x^{n_k}$. The leg count is the value
  at $x=1$, and the cardinality of the body is the value at $x=1$ of the
  formal derivative.
\item For finite sets with the cartesian structure, $\Iso(\cC)\iso(\mathbb N,\cdot)$
  and $\mathcal R(\cC)\iso\mathbb N[(\mathbb N,\cdot)]$.
\end{enumerate}
\end{example}

\section{Feline toposes}\label{sec:topos}

\begin{assumption}\label{ass:topos}
In this section $\cE$ is an elementary topos which has all small products. For
example, $\cE$ may be any Grothendieck topos.
\end{assumption}

Recall that an elementary topos has finite limits and finite colimits (in
particular an initial object), is cartesian closed, and has a subobject
classifier $\top\colon1\to\Omega$. We write $Y^X$ for exponentials in $\cE$.

\begin{proposition}\label{prop:exponential}
$\Fel(\cE)$ is cartesian closed. For cats $K=\langle X_a\rangle_{a\in A}$ and
$K'=\langle Y_b\rangle_{b\in B}$ the exponential is
\[
K'^{K}=\Bigl\langle\prod_{a\in A}Y_{\varphi(a)}^{X_a}\Bigr\rangle_{\varphi\in B^A},
\]
with evaluation $K'^K\times K\to K'$ of gait $(\varphi,a)\mapsto\varphi(a)$ and
components the composites
$\prod_{a'}Y_{\varphi(a')}^{X_{a'}}\times X_a\to Y_{\varphi(a)}^{X_a}\times X_a\to Y_{\varphi(a)}$.
\end{proposition}

\begin{proof}
$\Fel(\cE)$ has finite products by \cref{thm:limits}. For a cat
$M=\langle Z_q\rangle_{q\in Q}$, using \cref{cor:Fel-hom,cor:limit-formulas},
\begin{align*}
\Fel(\cE)(M\times K,K')&\iso\prod_{q\in Q}\prod_{a\in A}\coprod_{b\in B}\cE(Z_q\times X_a,Y_b)
\iso\prod_{q}\prod_{a}\coprod_{b}\cE(Z_q,Y_b^{X_a})\\
&\iso\prod_{q}\coprod_{\varphi\in B^A}\prod_{a}\cE(Z_q,Y_{\varphi(a)}^{X_a})
\iso\prod_q\coprod_{\varphi}\cE\Bigl(Z_q,\prod_aY_{\varphi(a)}^{X_a}\Bigr)\\
&\iso\Fel(\cE)(M,K'^K),
\end{align*}
naturally in $M$, where the third bijection is distributivity in $\Set$. Tracing
the identity of $K'^K$ through these bijections shows that the counit is the
stated evaluation.
\end{proof}

\begin{proposition}\label{prop:classifier}
Let $\Omega_{\Fel}\defeq\langle\Omega,1\rangle$ be the two-legged cat with leg
set $\{\mathsf{in},\mathsf{out}\}$, legs $\Omega$ at $\mathsf{in}$ and $1$ at
$\mathsf{out}$, and let $\top_{\Fel}\colon\Caudal\to\Omega_{\Fel}$ have gait
$*\mapsto\mathsf{in}$ and component $\top\colon1\to\Omega$. Then
$\top_\Fel$ is a subobject classifier of $\Fel(\cE)$.
\end{proposition}

\begin{proof}
By \cref{prop:mono} (which applies since $\cE$ has an initial object), a
monomorphism into $K=\langle X_a\rangle_A$ is a morphism $(\varphi,m)$ with
$\varphi$ injective and each $m_c$ monic. Two monomorphisms represent the same
subobject if and only if they have the same image $A'\subseteq A$ of gaits and,
for each $a\in A'$, the same subobject of $X_a$. Thus a subobject of $K$ is a
pair $(A',(S_a)_{a\in A'})$ with $A'\subseteq A$ and $S_a$ a subobject of $X_a$.
Given such a pair, define $\chi\colon K\to\Omega_\Fel$ with gait
$a\mapsto\mathsf{in}$ for $a\in A'$ and $a\mapsto\mathsf{out}$ otherwise, and with
component the classifying map $\chi_{S_a}\colon X_a\to\Omega$ for $a\in A'$ and
the unique map $X_a\to1$ otherwise. By \cref{cor:limit-formulas}, the
pullback of $\top_\Fel$ along $\chi$ has legs indexed by
$\{a:\chi(a)=\mathsf{in}\}=A'$, with leg $X_a\times_\Omega1=S_a$ at $a$; so
$\chi$ classifies $(A',S)$. Conversely, if $\chi'$ classifies $(A',S)$, then
the pullback computation shows that the gait of $\chi'$ sends exactly $A'$ to
$\mathsf{in}$, and that for $a\in A'$ the component of $\chi'$ classifies $S_a$;
components at legs sent to $\mathsf{out}$ are unique. So $\chi'=\chi$.
\end{proof}

\begin{theorem}[Feline toposes]\label{thm:topos}
If $\cE$ is an elementary topos with small products, then $\Fel(\cE)$ is an
elementary topos, whose subobject classifier is the two-legged cat
$\langle\Omega,1\rangle$.
\end{theorem}

\begin{proof}
$\Fel(\cE)$ has finite limits by \cref{thm:limits}, is cartesian closed by
\cref{prop:exponential}, and has a subobject classifier by
\cref{prop:classifier}.
\end{proof}

\begin{example}\label{ex:sierpinski}
For $\cE=\Set$, $\Fel(\Set)\eqv\Set^{\to}$ (\cref{cor:Fel-Set}), the topos of
sheaves on the Sierpi\'nski space. Under this equivalence the classifier
$\langle2,1\rangle$ corresponds to the function from a three-element set to a
two-element set which sends two elements to $\mathsf{in}$ and one to
$\mathsf{out}$, which is the familiar subobject classifier of $\Set^{\to}$.
\end{example}

\begin{remark}\label{rem:gluing}
If $\cE$ is a Grothendieck topos, extensivity of $\cE$ shows that $\Fel(\cE)$ is
equivalent to the comma category $\cE\downarrow\Delta$, where
$\Delta\colon\Set\to\cE$ sends a set $S$ to the copower $S\cdot1$: a cat
corresponds to the morphism $\body K\to\legs(K)\cdot1$. Since $\Delta$ is the
inverse image of the global sections geometric morphism, it preserves finite
limits, and $\cE\downarrow\Delta$ is the Artin gluing of $\Set$ and $\cE$ along
$\Delta$, which is a topos \cite{Johnstone2002}. We have preferred the direct
proof of \cref{thm:topos}, which uses no extensivity and applies to every
elementary topos with small products.
\end{remark}
